\subsection{Pure targets whose domains still require extraction}

The next three terminal proofs are mathematically complete in the existing
appendices and numbered proof package.  They are not yet suitable as
standalone formalization inputs, because their feasible domains are still
specified through geometric normal forms, support choices, or placement
suprema.  The correct revision is to expose that incompleteness rather than
hide geometry inside the word ``admissible.''

\begin{definition}[Problem S2-T3: four-label endpoint target]
\label{prob:s2-t3-endpoint}
For a real endpoint tuple
\[
 z=(x_-,x_+,c_-,c_+)\in[0,1]^4,
\]
define
\[
 \Psi_{\rm T3}(z)=F(c_-,x_-)+F(c_+,x_+)-1.
\]
The target optimization statement is
\[
 \Psi_{\rm T3}(z)<0
 \qquad(z\in\Omega_{\rm T3}),
\]
where $\Omega_{\rm T3}$ must be supplied as a finite union of explicit
real-variable contact-label domains.  The current exact proof determines this
union implicitly through the four-label audit, but the full list of
inequalities has not yet been consolidated into one formalization file.
\end{definition}

\begin{definition}[Problem S2-SC: adjacent-rescuer target]
\label{prob:s2-rescuer}
For
\[
 0\le b\le1,\qquad
 0\le c\le\frac12,\qquad
 \frac12\le u<1,
\]
put
\[
 E_{\rm sc}(c)=\frac{c+\sqrt{c^2-8c+4}}2,
 \qquad
 \theta(b,u)=\frac{b}{b+1-u},
\]
and
\[
 \Psi_{\rm SC}(b,c,u)
 =
 \max\{b,\theta(b,u)\}+E_{\rm sc}(c)-1.
\]
The target is
\[
 \Psi_{\rm SC}(b,c,u)\le0
 \qquad((b,c,u)\in\Omega_{\rm SC}).
\]
The exact rational local profile proves this statement, but the present
manuscript still defines $\Omega_{\rm SC}$ by a translated-role normal form.
A formalization-ready revision must replace that normal form by its explicit
polynomial and sign constraints.
\end{definition}

\begin{definition}[Problem S2-VD: Vd radial targets]
\label{prob:s2-vd-terminal}
The adjacent target is a finite maximum of the differences between each
vertex-side bridge endpoint and the corresponding center-side start.  The
nonadjacent target is a finite maximum of
\[
 s-a_{\rm R},\qquad
 s-a_{\rm L},\qquad
 d_\tau-\min\{a_{\rm R},a_{\rm L}\},\qquad
 C_\tau-1+\min\{a_{\rm R},a_{\rm L}\}.
\]
Each maximum must be strictly negative on its real feasible domain
$\Omega_{\rm VD}^{\rm adj}$ or $\Omega_{\rm VD}^{\rm non}$.  The quarter-radial
and Vd-corner lemmas prove these signs, but the domains still contain
placement-defined suprema.  Their elimination into polynomial inequalities
is a remaining interface task.
\end{definition}

The Vd1--supercritical two-chart replacement is intentionally absent from
this list.  It changes the geometric roles while preserving skeleton
coverage; it is a geometric reduction rather than a scalar optimization.

\subsection{Formalization status}

\begin{table}[H]
\centering
\small
\renewcommand{\arraystretch}{1.12}
\begin{tabular}{@{}p{.15\textwidth}p{.27\textwidth}p{.19\textwidth}p{.28\textwidth}@{}}
\toprule
identifier&input and objective&interface status&remaining work\\
\midrule
S2-E1&three variables, two evaluations of $F$&ready&
branch split, radical sign checks, strict equality exclusion\\
S2-E2&three variables, two evaluations of $F$&ready&
same\\
S2-R1&three variables, three-step recurrence&ready&
encode the CE1 branch audit and chord certificates\\
S2-R2&three variables, three-step recurrence&ready&
encode the at-least-one-threshold disjunction\\
S2-T3&four endpoint variables, two evaluations of $F$&partial&
write the complete finite union $\Omega_{\rm T3}$ explicitly\\
S2-SC&three variables, one envelope objective&partial&
eliminate the translated-role normal form\\
S2-VD&adjacent/nonadjacent radial objectives&partial&
eliminate placement suprema and list branch domains\\
replacement&no scalar objective&not applicable&
formalize the two geometric chart constructions directly\\
\bottomrule
\end{tabular}
\caption{Formalization readiness of the Strategy~2 terminal calculations.
``Partial'' refers only to the geometry-to-domain interface; the corresponding
mathematical terminal theorem retains its recorded proved status.}
\label{tab:s2-formalization-status}
\end{table}

The recommended formalization order is S2-E1, S2-E2, S2-R2, S2-R1, S2-SC,
S2-VD, and finally S2-T3.  The first two provide the smallest test of the
piecewise function and strict-branch certificate; S2-T3 should be last because
its feasible set has the largest contact-label expansion.
